\subsection{Escenario de Pruebas}
El prototipo será probado en la Unidad Profesional Interdisciplinaria de
Biotecnología IPN (UPIBI) en Planta Piloto en el Laboratorio de Biorreactores, en un
entorno con condiciones de higiene ideales para realizar fermentaciones de kéfir de
leche. Para validar el funcionamiento del sistema de monitoreo propuesto, se
plantea un escenario de pruebas controlado utilizando, cuatro frascos, cada uno con
diferentes condiciones de fermentación. De manera que 2 frascos serán preparados con 2 cepas de búlgaros totalmente diferentes a través de 3 usuarios iniciales y el
3er frasco será preparado con una cepa de búlgaros derivada de 2 productores distintos
, siendo este el usuario raíz.
Se proponen dos escenarios de pruebas:
En el primero se prepararán con diferentes parámetros distintos (cantidad de
búlgaros y tipo de leche) 3 fermentaciones ideales (con temperatura ideal) y una no
ideal (con una temperatura ideal) que estén próximas a llegar al rango de pH (4.2-
4.6), verificando de forma individual el desarrollo de cada fermentación a través de
sus respectivos usuarios mediante las páginas web y posterior comparar todas las
fermentaciones en curso a través del usuario administrador. Esto con el objetivo de
poder observar el desarrollo de cada fermentación y visualizar las alarmas que emite
sobre cada uno.
Para el segundo escenario se propone que posterior a validar las anteriores
fermentaciones iniciar 4 fermentaciones, 3 de manera ideal, es decir usando un tipo
de leche que su pH este entre 6.6-6.8 y una temperatura inicial optima, y la 4ta
fermentación se iniciara de manera no optima, es decir que la temperatura inicial y
el pH no sean óptimos. De manera que se verificara de forma individual cada
fermentación a través de sus respectivos usuarios mediante la página web y
posterior comparar todas las fermentaciones en curso a través del usuario
administrador. Esto con el objetivo de poder visualizar que al momento de detectar
parámetros inicales no ideales el sistema no inicia el proceso de monitoreo por parte
de los sensores en la fermentación.